\documentclass[10pt,a4paper]{amsart}
\usepackage[margin=19mm]{geometry}
\usepackage{amsmath,amssymb,mathtools}
\usepackage[scr=rsfs,cal=euler]{mathalfa}
\usepackage{xfrac}
\usepackage[unicode,hidelinks,bookmarks=false]{hyperref}
\newcommand{\ie}{i.e.\ }
\newcommand{\cf}{cf.\ }
\newcommand{\resp}{respectively\ }
\newcommand{\Subsection}{\subsection}
\providecommand{\nobreakdash}{\nobreak\hbox{-}}
\setlength{\emergencystretch}{2em}
\multlinegap0pt
\usepackage{mathtools}
\usepackage{amssymb}
\usepackage[shortlabels]{enumitem}
\usepackage{dsfont}
\usepackage{xcolor}
\newtheorem{theorem}{Theorem}
\newtheorem{proposition}{Proposition}
\newtheorem{corollary}{Corollary}
\newcommand{\PSH}{{\rm PSH}}
\title{Singularities vs non-pluripolar Monge--Amp\`{e}re masses}
\author{Quang-Tuan Dang, Hoang-Son Do, Hoang Hiep Pham}
\date{Source: arXiv:2503.07534v1 (CC BY 4.0)}
\begin{document}
\maketitle

\noindent\emph{Source and licence.} Singularities vs non-pluripolar Monge--Amp\`{e}re masses. Version arXiv:2503.07534v1 (CC BY 4.0), distributed by arXiv under the Creative Commons Attribution 4.0 International licence, http://creativecommons.org/licenses/by/4.0/, as recorded in the arXiv licence metadata for this exact version and shown on the abstract page https://arxiv.org/abs/2503.07534v1. Full licence notice: \texttt{licenses/arxiv-2503.07534-CC-BY-4.0.txt}.\par\smallskip

\noindent\textit{Editorial scope.} Counted unit: the article's corollary corollary (source0.tex lines 749--760). The article's complete original proof of it is delivered with it (source0.tex lines 762--773, one \texttt{proof} environment of 12 lines). No mathematics is written, repaired or re-derived: the statement and the proof are the article's own lines, in the article's own notation, and no retained line is substituted or re-flowed.  Retained uncounted as the source-local context the proof needs: lines 307--316; lines 443--445.  Nothing was omitted from any retained span, so the package carries no omission record.  No retained line contains a citation, so the wrapper carries no bibliography and no bibliography entry.  No retained line contains a figure environment, an image-inclusion command or drawing code, so no image is delivered and the package declares no asset dependency.  Editorial formatting only, in addition: an AMS \texttt{amsart} preamble that restates the article's own declarations and macros used by the retained lines, namely the environments \texttt{theorem}, \texttt{proposition}, \texttt{corollary} on separate counters, together with the standard packages those lines need.  The retained environments print in this document as Theorem 1, Proposition 1, Corollary 1 in order of appearance, whereas the article prints the same items with the numbering of its own full document.  The complete native source of the article as distributed in the arXiv e-print for this exact version is bundled unchanged as \texttt{sources/174515/source0.tex}.
\begin{theorem}\label{thm: thm1}
	%Let $\varphi,\psi\in\PSH(X,\theta)$ with positive masses. If $\int_X\theta_\varphi^n=\int_X\theta_\psi^n>0$, 
Let $(X, \omega)$ be a compact K\"{a}hler manifold and $\theta$ a smooth closed real (1, 1)-form on $X$ whose cohomology class is big.
Let $\phi\in\PSH(X,\theta)$. If $\varphi,\psi\in\mathcal{E}(X,\theta,\phi)$, then for each $c>0$, there exists a function $ h\in\mathcal{E}(X,\omega)$ such that $$|\varphi-\psi|\leq -ch.$$
Conversely, if $\varphi,\psi\in\PSH(X,\theta)$ such that $|\varphi-\psi|\leq -ch$ for some $c>0$ and $h\in\mathcal{E}(X,\omega)$, then both $\varphi$ and $\psi$ belong to $\mathcal{E}(X,\theta,\max(\varphi,\psi))$.

In particular, $\varphi\in\mathcal{E}(X,\theta,\phi)$ if and only if $\varphi$ is more singular than $\phi$ and $\varphi\geq \phi+ch $ for some $c>0$ and $h\in\mathcal{E}(X,\omega)$.
% $$\int_X\theta_\varphi^n=\int_X\theta_\psi^n.$$
% Furthermore, if $\theta_\varphi^n=\theta_\psi^n$ then $u-v$ is constant. 
\end{theorem}  

\begin{proposition}\label{prop: homega}
Let $\{\theta\}$ be a big cohomology class and $\phi\in\PSH(X,\theta)$. Let $\varphi\in\mathcal{E}(X,\theta,\phi)$ and $\psi\in \PSH(X,\theta,\phi)$. Then for any $b>0$, $P_\omega(b\varphi-b\psi)\in\mathcal{E}(X,\omega)$.
\end{proposition}

\begin{corollary}
 Let $\{\eta\}$ and $\{\theta\}$ be big cohomology classes. Let $\phi\in\PSH(X,\theta)$. Then the following are equivalent.
\hfill
\begin{enumerate}
\item $\PSH(X,\eta)\cap\mathcal{E}(X,\theta,\phi)\neq \varnothing$.

\item There exists $u\in\PSH(X,\eta)$, $v\in\mathcal{E}(X,\theta,\phi)$, $c>0$, and $h\in\mathcal{E}(X,\omega)$ such that \[u\geq v+ ch. \]

\item $\mathcal{E}(X,\eta)\cap\PSH(X,\theta,\phi)\subset \mathcal{E}(X,\theta,\phi)$.
\end{enumerate}

\end{corollary}

\begin{proof}
 We see that $(1)\Rightarrow (2)$ is trivial. Indeed, we can take $u=v\in \PSH(X,\eta)\cap\mathcal{E}(X,\theta,\phi)$ and $h=0$.
% Assume $u\in \PSH(X,\eta)\cap\mathcal{E}(X,\theta,\phi)$. Thanks to Proposition~\ref{prop: homega}, we obtain $h\colonequals P_\omega(c^{-1}(u-\phi))\in\mathcal{E}(X,\omega)$. Hence,
% \[u\geq \phi+ch. \]

For $(2)\Rightarrow (3)$, let $\varphi\in \mathcal{E}(X,\eta)\cap \PSH(X,\theta,\phi)$. Thanks to Proposition \ref{prop: homega}, we have $h'\coloneqq P_\omega(\varphi-V_\eta)\in\mathcal{E}(X,\omega)$. By assumption, there exists $u\in\PSH(X,\eta)$, $v\in\mathcal{E}(X,\theta,\phi)$, $c>0$, and $h\in\mathcal{E}(X,\omega)$ such that \[u\geq v+ ch. \]
 It follows that 
 \[\varphi\geq V_\eta+h'\geq u-\sup_X u+h'\geq v+ch+h'-\sup_X u. \]
 Theorem \ref{thm: thm1} yields $\int_X\theta_\varphi^n=\int_X\theta^n_v=\int_X\theta_\phi^n$, hence $\varphi\in\mathcal{E}(X,\theta,\phi)$.

 For $(3)\Rightarrow (1)$, if $\mathcal{E}(X,\eta)\cap\PSH(X,\theta,\phi)=\varnothing$, we are done. Otherwise, $u\in \mathcal{E}(X,\eta)\cap\PSH(X,\theta,\phi)$, so $u\in\mathcal{E}(X,\theta,\phi)$. 
\end{proof}

\end{document}
